\chapter{\ifenglish Background Knowledge and Theory\else ทฤษฎีที่เกี่ยวข้อง\fi}

การพัฒนาเกมคอมพิวเตอร์จำเป็นต้องอาศัยความรู้จากหลายด้าน เช่น การเขียนโปรแกรม การออกแบบระบบซอฟต์แวร์ การออกแบบกราฟิก และการออกแบบประสบการณ์ผู้ใช้ โดยบทนี้จะอธิบายทฤษฎีและแนวคิดพื้นฐานที่เกี่ยวข้องกับการพัฒนาเกมในโครงงาน เพื่อให้ผู้อ่านสามารถเข้าใจแนวทางการออกแบบและพัฒนาระบบของเกมได้ง่ายขึ้น

\section{Game Development}

Game Development คือกระบวนการสร้างเกมคอมพิวเตอร์ตั้งแต่ขั้นตอนการออกแบบแนวคิดของเกม (Game Concept) การออกแบบระบบการเล่น (Gameplay Design) การพัฒนาโปรแกรม (Programming) ไปจนถึงการทดสอบและปรับปรุงเกม (Testing and Debugging)

โดยทั่วไปกระบวนการพัฒนาเกมจะประกอบด้วยขั้นตอนสำคัญ ได้แก่

\begin{itemize}
\item การออกแบบแนวคิดของเกม (Game Design)
\item การออกแบบระบบการเล่น (Gameplay Mechanics)
\item การพัฒนาโปรแกรม (Game Programming)
\item การสร้างกราฟิกและแอนิเมชัน (Art and Animation)
\item การทดสอบระบบ (Game Testing)
\end{itemize}

โครงงานนี้ใช้กระบวนการพัฒนาเกมในลักษณะของการสร้างต้นแบบ (Prototype) เพื่อทดสอบแนวคิดของระบบเกมก่อนการพัฒนาเต็มรูปแบบ

\section{Unity Game Engine}

Unity เป็น Game Engine ที่ได้รับความนิยมอย่างแพร่หลายในการพัฒนาเกมทั้งแบบ 2 มิติและ 3 มิติ โดยรองรับการพัฒนาเกมสำหรับหลายแพลตฟอร์ม เช่น Windows, macOS, Android และ iOS

Unity มีคุณสมบัติสำคัญที่ช่วยให้การพัฒนาเกมมีประสิทธิภาพ เช่น

\begin{itemize}
\item ระบบ Physics สำหรับการจำลองการเคลื่อนไหวของวัตถุ
\item ระบบ Animation สำหรับควบคุมการเคลื่อนไหวของตัวละคร
\item ระบบ Scene Management สำหรับจัดการฉากในเกม
\item ระบบ Script ที่ใช้ภาษา C\# ในการพัฒนา
\end{itemize}

ในโครงงานนี้ Unity ถูกใช้เป็นเครื่องมือหลักในการพัฒนาเกม เนื่องจากมีเครื่องมือที่ช่วยให้สามารถพัฒนาเกมต้นแบบได้อย่างรวดเร็ว

\subsection{C\# Programming Language}

C\# เป็นภาษาโปรแกรมเชิงวัตถุ (Object-Oriented Programming Language) ที่ถูกใช้ใน Unity สำหรับเขียนสคริปต์ควบคุมการทำงานของเกม เช่น การควบคุมตัวละคร การจัดการศัตรู และระบบต่าง ๆ ของเกม

คุณสมบัติสำคัญของภาษา C\# ได้แก่

\begin{itemize}
\item รองรับการเขียนโปรแกรมเชิงวัตถุ
\item มีโครงสร้างที่ชัดเจนและอ่านง่าย
\item รองรับการพัฒนาโปรแกรมขนาดใหญ่
\end{itemize}

\subsubsection{Game Script}

Game Script คือโปรแกรมที่ใช้ควบคุมพฤติกรรมของวัตถุภายในเกม เช่น การเคลื่อนไหวของตัวละคร การโจมตีของศัตรู หรือการทำงานของระบบเกมต่าง ๆ

ใน Unity สคริปต์มักจะถูกผูกเข้ากับ Game Object ภายใน Scene เพื่อกำหนดพฤติกรรมของวัตถุนั้น ๆ

\subsubsection{Component System}

Unity ใช้โครงสร้างแบบ Component-Based Architecture โดยวัตถุในเกมจะเรียกว่า Game Object และสามารถเพิ่ม Component ต่าง ๆ เข้าไปเพื่อกำหนดคุณสมบัติ เช่น

\begin{itemize}
\item Transform Component
\item Rigidbody
\item Collider
\item Animator
\end{itemize}

โครงสร้างดังกล่าวช่วยให้สามารถพัฒนาระบบเกมได้อย่างยืดหยุ่นและง่ายต่อการขยายระบบในอนาคต

\section{2D Game System}

เกมในโครงงานนี้เป็นเกมแบบสองมิติ (2D Game) ซึ่งใช้ระบบ Sprite และ Tilemap สำหรับการสร้างฉากและตัวละครภายในเกม

ระบบหลักที่ใช้ในการพัฒนาเกม 2 มิติ ได้แก่

\begin{itemize}
\item Sprite Rendering สำหรับการแสดงผลภาพตัวละคร
\item Animation System สำหรับการเคลื่อนไหวของตัวละคร
\item Physics System สำหรับการตรวจจับการชน
\item Input System สำหรับรับคำสั่งจากผู้เล่น
\end{itemize}

\section{Gameplay Mechanics}

Gameplay Mechanics คือกฎและระบบที่กำหนดรูปแบบการเล่นของเกม ซึ่งมีผลต่อประสบการณ์ของผู้เล่นโดยตรง

ในโครงงานนี้ ระบบ Gameplay หลักประกอบด้วย

\begin{itemize}
\item ระบบการสำรวจแผนที่ (Exploration)
\item ระบบการต่อสู้ (Combat System)
\item ระบบไอเทมและอุปกรณ์ (Item System)
\item ระบบมินิเกม (Mini-game System)
\item ระบบบอส (Boss System)
\end{itemize}

\section{\ifenglish%
\ifcpe CPE \else ISNE \fi knowledge used, applied, or integrated in this project
\else%
ความรู้ตามหลักสูตรซึ่งถูกนำมาใช้หรือบูรณาการในโครงงาน
\fi
}

ในการพัฒนาโครงงานนี้ มีการนำความรู้จากหลายวิชาที่ได้ศึกษาในหลักสูตรมาใช้ เช่น

\begin{itemize}
\item วิชาการเขียนโปรแกรมคอมพิวเตอร์
\item วิชาโครงสร้างข้อมูลและอัลกอริทึม
\item วิชาวิศวกรรมซอฟต์แวร์
\item วิชาการออกแบบระบบคอมพิวเตอร์
\end{itemize}

ความรู้จากวิชาเหล่านี้ถูกนำมาใช้ในการออกแบบโครงสร้างโปรแกรม การจัดการข้อมูล และการพัฒนาระบบต่าง ๆ ภายในเกม

\section{\ifenglish%
Extracurricular knowledge used, applied, or integrated in this project
\else%
ความรู้นอกหลักสูตรซึ่งถูกนำมาใช้หรือบูรณาการในโครงงาน
\fi
}

นอกจากความรู้ตามหลักสูตรแล้ว โครงงานนี้ยังมีการศึกษาและเรียนรู้เพิ่มเติมจากแหล่งข้อมูลภายนอก เช่น เอกสารออนไลน์ บทความเกี่ยวกับการพัฒนาเกม และบทเรียนจากชุมชนนักพัฒนาเกม

ความรู้ที่ได้จากการศึกษาด้วยตนเอง เช่น การใช้งาน Unity Game Engine การออกแบบระบบเกม และเทคนิคการพัฒนาเกม ถูกนำมาประยุกต์ใช้ในการพัฒนาโครงงานนี้เพื่อให้สามารถสร้างเกมต้นแบบที่มีระบบการเล่นสมบูรณ์มากยิ่งขึ้น